\documentclass[11pt]{article}

% ---------- Layout ----------
\usepackage[letterpaper,margin=1in]{geometry}
\usepackage{setspace}
\setstretch{1.05}

% ---------- Math & text ----------
\usepackage[T1]{fontenc}
\usepackage[utf8]{inputenc}
\usepackage{amsmath,amssymb}
\usepackage{microtype}
\usepackage{longtable}
\usepackage{booktabs}
\usepackage{pdflscape} % Required for wide appendix-style tables when needed

% ---------- Figures & tables ----------
\usepackage{graphicx}
\usepackage{float}
\usepackage{array}
\usepackage{caption}

% ---------- Links ----------
\usepackage[colorlinks=true,linkcolor=blue,citecolor=blue,urlcolor=blue]{hyperref}

% ---------- Bibliography ----------
\usepackage[authoryear]{natbib}

% ---------- Local commands ----------
\newcommand{\projecttranslation}{\textsc{[Project translation]}}
\newcommand{\sourcegap}{\textsc{[Source gap]}}
\newcommand{\SHT}{Hirano--Stiglitz--Toda}

% ---------- Title block ----------
\title{Land, Credit, Overvaluation, and Debt Rollover:\
A Systematic Narrative Review of the Hirano--Stiglitz--Toda Branch\thanks{This document is a working synthesis prepared for the Fondecyt project on Latin American land-rent financialization. It benefited from AI-assisted editing for syntax refinement, structural suggestions, and stylistic consistency checks. All analytical arguments, historical interpretations, and substantive content remain the authors' sole responsibility.}\\
\begin{large}
\textbf{WORK IN PROGRESS / PLEASE DO NOT CIRCULATE}
\end{large}
}
\author{Diego Polanco\thanks{Department of Economics, University of Massachusetts Amherst\\
\texttt{dpolanconeco@umass.edu}} \and Ivo Gasic}
\date{\today}

\begin{document}
\maketitle

\begin{abstract}
This review reconstructs the bounded Hirano--Stiglitz--Toda branch on land, credit, overvaluation, and debt rollover. The central claim is that land is not a neutral background asset: because it is non-produced, durable, scarce, and able to function as a production factor, store of value, and collateral object, land can reorganize macro-financial relations. Land-price inflation can raise measured wealth without expanding productive capacity; credit routed toward land and real estate can enlarge balance sheets while reducing long-run growth; unbalanced growth can separate land prices from rent fundamentals; and land can modify the relation between returns, growth, and debt rollover. The review identifies four mechanisms---capital-formation crowding-out, sectoral credit and collateral allocation, unbalanced-growth land overvaluation, and the land-mediated critique of simple \(R\)-versus-\(G\) debt-rollover reasoning---and evaluates what can be carried into a Fondecyt project on Latin American land-rent financialization. Its conclusion is deliberately bounded: the SHT corpus provides a formal mechanism set, not a ready-made account of Latin American urban political economy. Translation requires additional institutional, spatial, and regional evidence.
\end{abstract}

\vspace{1cm}

\noindent\textbf{Keywords:} Land Rent; Land Speculation; Financialization; Credit Allocation; Land Overvaluation; Debt Rollover; Unbalanced Growth; Productive Capital Formation.

\noindent\textbf{JEL Codes:} B51; E22; E44; G12; O11; P16; R30.

\pagebreak

\section*{Review question}

How does the Hirano--Stiglitz--Toda literature explain the relationship between land-price inflation, credit expansion, unbalanced growth, productive capital formation, land overvaluation, and debt rollover, and what can be carried into a Fondecyt project on Latin American land-rent financialization?

\section*{Executive synthesis}

The bounded \SHT{} branch does not treat land as a neutral background asset. Land is non-produced, durable, scarce, and able to function as production factor, store of value, and collateral object. That combination creates the macro-financial problem. Land-price inflation can raise measured wealth without adding productive capacity. Credit routed toward land or real estate can enlarge balance sheets while reducing long-run growth. Differential productivity growth can pull land prices away from land rents. Once land enters the asset structure, the relation between returns, growth, and debt rollover also changes.

The core argument is not the generic claim that financialization raises land rent. The sharper claim is that land valuation becomes macroeconomically central when credit, collateral, savings, and growth organize around a non-produced asset. Four mechanisms carry the review: land speculation crowds out capital formation; real-estate credit and collateral amplify land-price inflation; unbalanced growth can separate land prices from rent fundamentals; and land can modify \(R\) versus \(G\) and the possibility of infinite debt rollover. These mechanisms are strong as formal theory and thin as institutional description. For the Fondecyt project, they must be reconstructed first on their own terms, then translated into Latin American urban land-rent and financialization contexts.

\section{Introduction: the object of the review}

\paragraph{P01.} Land-price inflation matters because it can raise balance-sheet wealth without expanding productive capacity. Land is not produced like machines, infrastructure, or structures. The same price increase can appear as private wealth, collateral value, or capital gain while leaving the produced capital stock largely unchanged. Stiglitz's wealth-measurement papers provide the conceptual base: measured wealth can rise through capitalized rents, including land and monopoly rents, even when real capital accumulation and national saving do not rise in proportion. The problem is therefore not only mismeasurement. It is the possibility that an economy registers asset-value gains while its stock of produced means of production grows slowly, stagnates, or becomes relatively smaller within the balance sheet \citep[pp.~1--2, 4, 11]{stiglitz2015MeasurementWealth}; \citep[pp.~10--11, 28]{stiglitz2015WealthResidual}.

\paragraph{P02.} The organizing unit is the relation, not the paper. Land versus productive capital defines the object. Credit versus land price defines the financial channel. Unbalanced growth versus rent fundamentals defines overvaluation. Land, \(R\), \(G\), and debt rollover define the debt-sustainability channel. Wobbly dynamics remains supplementary: it explains endogenous instability and phase transitions, but it does not displace the four central mechanisms \citep[pp.~7, 24, 46--47]{hiranoStiglitz2022LandSpeculation}; \citep[pp.~1--2, 10]{hiranoStiglitz2022WobblyEconomy}.

\paragraph{P03.} Reconstruction comes before critique. The \SHT{} branch is treated first as a formal macro-financial research program with its own objects: land, land rent, land price, productive capital, wealth, collateral, leverage, \(R\), \(G\), and debt rollover. Its methods are overlapping-generations models, rational expectations, collateral constraints, no-arbitrage conditions, and unbalanced growth. The Fondecyt translation comes later. These models are not Latin American urban political economy; they identify mechanisms that can travel only after their institutional assumptions are made explicit.

\section{Method and corpus}

\paragraph{P04.} The corpus contains fourteen registered sources and three evidence-control layers: ninety-eight per-paper question notes, thirty-two cluster synthesis notes, and a claim-level extraction matrix. The question notes preserve paper-level provenance; the cluster notes identify cross-paper relations; the matrix separates definitions, mechanisms, causal chains, variables or parameters, and theory/evidence status. These layers are infrastructure, not citable sources. Claims enter this draft only when they can be routed back to registered source papers. Where the notes or matrix support a claim but do not supply publication-ready citation detail, the draft marks the problem rather than inventing a page or reference.

\paragraph{P05.} The extraction protocol keeps five dimensions separate. Definitions identify the objects: land, land rent, housing rent, real-estate value, productive capital, wealth, bubble, and overvaluation. Mechanisms explain how those objects connect. Causality orders the movement, for example from credit easing to land demand, land prices, and capital crowding-out. Variables and parameters show where the model does its work: prices and rents \((P_t,D_t)\), capital and investment \((K_t,k_{t+1})\), returns and growth \((R,G,r)\), credit conditions \((\theta,\theta^x)\), leverage, and substitution conditions \((\sigma)\). Theory/evidence status then separates formal results, empirical motivations, and policy implications.

\paragraph{P06.} Cluster synthesis supplies infrastructure, not finished prose. C1 grounds the wealth-capital distinction. C2 organizes wobbly dynamics and instability. C3 gathers credit, low rates, and crowding-out. C4 links unbalanced growth, leverage, land bubbles, and overvaluation. C5 handles land, \(R\) versus \(G\), and debt rollover. C6 is corpus-level synthesis rather than a bounded cluster, so its claims require extra caution. It can organize a synthesis or flag a transportability problem, but it cannot substitute for source-level support from the papers.

\section{Conceptual foundations: land, rent, wealth, and productive capital}

\paragraph{P07.} Productive capital and wealth are not the same object. Productive capital consists of produced assets used in production. Wealth is wider: ownership claims over productive capital, land, capitalized rents, monopoly rents, intellectual-property rents, and other claims. Rising wealth is therefore not automatically capital deepening. In Stiglitz's formulation, the wealth residual is the gap between observed increases in wealth and the smaller increase implied by real capital accumulation and national saving. The residual matters because it identifies balance-sheet expansion whose source is not new produced capital but the capitalization of rent claims and asset positions \citep[pp.~10--11, 28]{stiglitz2015WealthResidual}; \citep[pp.~1--2, 4, 11]{stiglitz2015MeasurementWealth}.

\paragraph{P08.} Land rent and land price must also be kept apart. Land rent is a flow or dividend generated by land. Land price is the capitalized asset price of the land. The Land Overvaluation Theorem turns on the case where they diverge. The fundamental value of land is the present value of expected future land rents under the model's discounting structure. Land is overvalued when its market price contains a component above that fundamental value. This is not the claim that every high land price is a bubble. It is a specific relation between an asset price, the rent stream that supposedly anchors it, and the discounting rule that prices that stream \citep[pp.~1, 3--4, 13]{hiranoToda2025UnbalancedGrowth}.

\paragraph{P09.} Housing and real estate are composite objects. Real estate joins land to produced structures. Housing rent includes the service flow from structures as well as the land component. Pure land rent is the return to the non-produced site. Toda's short paper blocks a misleading inference from observed housing rent-yield ratios: a stationary aggregate housing rent-to-price ratio does not rule out a land bubble if the analyst has not separated the return on structures from the pure land dividend. Housing data can look anchored while the land component is the part undergoing overvaluation \citep[pp.~1--3]{toda2025LandBubbles}.

\paragraph{P10.} These distinctions discipline the analysis of financialization. A bank can expand real-estate or land-collateral credit, land prices can rise, and measured wealth and collateral values can jump immediately. None of this proves that the economy has produced more capital or that land rent has grown at the same pace as the asset price. The credit channel links the concepts: credit raises purchasing power for a fixed asset; the asset price raises collateral values; collateral values can support more credit. The SHT branch uses this distinction to show how an economy can look richer on balance sheets while becoming more exposed to lower productive investment, fragile collateral values, or unstable asset-price paths \citep[pp.~16--18]{stiglitz2015LandCredit}; \citep[pp.~1, 15, 33, 38]{hiranoStiglitz2024CreditLandGrowth}.

\section{Mechanism I: land-price inflation and capital-formation crowding-out}

\paragraph{P11.} Crowding-out begins when land competes with productive capital for savings. In the OLG version, young agents allocate finite savings between capital and land. A higher land price absorbs more of that savings envelope and leaves less for real capital accumulation. A simplified normalized expression for the mechanism is
\[
  k_{t+1}+P_t=s_t(w_t+e),
\]
where next-period capital plus the current land price must fit inside savings out of wages and endowment. In the general model, the exact budget and market-clearing conditions are richer; the substantive point is the same. Holding savings fixed, an increase in \(P_t\) is not an increase in produced capital. It is a larger claim on the same savings pool, leaving a smaller residual for \(k_{t+1}\). This is the sense in which land holdings can crowd out real capital accumulation \citep[pp.~7, 24, 46--47]{hiranoStiglitz2022LandSpeculation}.

\paragraph{P12.} The chain is more than accounting. A land boom raises the price of a fixed asset. Agents treat land as a store of value, and resale expectations sustain demand. Higher land prices reduce capital accumulation, lower future wages and output, and can raise the return on scarce productive capital. That higher return can require faster land-price appreciation for arbitrage to hold. The boom therefore carries its own instability: it needs rising prices to remain attractive, while the same price increase weakens the capital base that supports future income. Once the path hits its bound, expectations can reverse and the boom can endogenously become a crash \citep[pp.~7, 24, 46--47]{hiranoStiglitz2022LandSpeculation}.

\paragraph{P13.} The Henry George paper gives this mechanism a policy-theoretic form. Land speculation diverts scarce savings away from productive investment. A tax on land, or on returns to land speculation, can redirect resources toward productive capital and improve growth and welfare under the model's assumptions. The point is not only revenue. Land taxation changes portfolio incentives where an asset in fixed supply competes with produced capital for savings. It lowers the after-tax attractiveness of holding land for capital gains and can shift the savings margin back toward productive investment \citep[pp.~327, 329, 335, 345]{hiranoStiglitz2025HenryGeorge}.

\paragraph{P14.} The crowding-out mechanism is conditional. It does not say that every land purchase reduces growth. The result depends on land functioning as an alternative store of value, on savings or financing being limited at the relevant margin, on land-price appreciation being attractive relative to productive investment, and on land purchases not creating new productive capacity. Those conditions are plausible in financialized land markets, but they must be shown empirically. For Latin American cases, this is a \sourcegap{}: the current corpus does not itself provide evidence on credit composition, land tenure, informal housing production, or metropolitan land markets.

\section{Mechanism II: credit, collateral, leverage, and monetary policy}

\paragraph{P15.} Credit sharpens the crowding-out claim by shifting attention from quantity to allocation. Hirano and Stiglitz argue that the macroeconomic effect of credit expansion depends less on aggregate credit than on where the credit goes. Credit for productive capital can raise productivity and growth. Credit for real estate or land can fuel speculation, inflate asset prices, and reduce long-run growth. The long-run growth paper states the sectoral-credit hypothesis directly: it is ``not so much aggregate credit expansion'' that matters, but sectoral credit expansion. A larger credit stock may coincide with lower growth if the marginal loan finances land purchases, down payments, or real-estate collateral positions rather than capital formation in the productive sector \citep[pp.~1, 15, 33, 38]{hiranoStiglitz2024CreditLandGrowth}.

\paragraph{P16.} The low-interest-rate paper works through the same channel. Rapid credit expansion is associated with later growth slowdowns when credit flows disproportionately into real estate, while credit expansion toward manufacturing is associated with higher productivity. The model explains why: lower rates or looser collateral constraints reduce the down payment required to acquire land, raise the leveraged return on land speculation, and move portfolios away from productive capital. What looks like easier finance in partial equilibrium can become lower productive investment once land-price effects enter general equilibrium \citep[p.~1]{hiranoStiglitz2025LowInterestPolicy}; \citep[pp.~1, 15, 33, 38]{hiranoStiglitz2024CreditLandGrowth}.

\paragraph{P17.} The collateral distinction carries the mechanism. Hirano and Stiglitz separate collateral constraints tied to productive capital from those tied to land. Relaxing capital collateral can crowd in productive investment. Relaxing land or real-estate collateral can crowd out productive investment if it raises the payoff to land speculation. The effect depends on the relative pledgeability of capital and land, represented here as \(\theta\) for capital collateral and \(\theta^x\) for land collateral. A higher \(\theta\) expands borrowing capacity attached to productive capital; a higher \(\theta^x\) expands borrowing capacity attached to land. The signs can therefore diverge: loosening capital finance can raise \(k_{t+1}\), while loosening land finance can raise \(P_t\) and reduce the savings left for productive investment \citep[pp.~4, 7, 10, 15, 33, 38]{hiranoStiglitz2024CreditLandGrowth}; \citep[pp.~1--2, 16--17, 19, 29]{hiranoStiglitz2025LowInterestPolicy}.

\paragraph{P18.} Stiglitz's land-and-credit analysis supplies the bridge to the later Hirano--Stiglitz models. In the Part IV paper, credit availability and bank decisions can drive land prices upward, raising measured wealth without increasing the equilibrium capital stock. Land serves as collateral, and credit conditions can inflate that collateral value. The mechanism is circular: credit supports land prices, higher land prices support collateral values, and higher collateral values support more credit. Because the collateral object is non-produced land, the balance-sheet effect can precede, exceed, or substitute for productive accumulation \citep[pp.~16--18]{stiglitz2015LandCredit}.

\paragraph{P19.} The financial channel makes monetary policy ambiguous, not irrelevant. In a Tobin-style intuition, lower rates or higher liquidity push agents from money into productive capital. In the SHT branch, the asset menu includes land, and land can absorb liquidity that might otherwise fund capital formation. Expansionary policy can therefore support or retard growth depending on whether it relaxes productive investment finance or land-collateral finance. The result is a compositional test: which balance sheets receive leverage and which asset prices respond first \citep[pp.~1--4]{hiranoStiglitz2025LowInterestPolicy}; \citep[pp.~1, 15, 33, 38]{hiranoStiglitz2024CreditLandGrowth}.

\section{Mechanism III: unbalanced growth and the Land Overvaluation Theorem}

\paragraph{P20.} Unbalanced growth has a narrow meaning in this branch. It refers to differential productivity growth between land and non-land factors or sectors. In the land overvaluation model, labor or human-capital productivity grows faster than land productivity. If the elasticity of substitution between land and labor exceeds one, output can keep expanding while land's productive contribution declines. The condition is not simply uneven development; it is a formal relation among factor productivity, substitutability, and the rent stream generated by land. Land rents grow slowly relative to wages and aggregate income, while land remains a durable store of value purchased out of rising savings \citep[pp.~1, 3--4, 13]{hiranoToda2025UnbalancedGrowth}.

\paragraph{P21.} The theorem's plain-language result is stark: land can become systematically overvalued even as its productive role declines. Theorem 1 requires \(\sigma>1\) and a convergent expected infinite sum involving the ratio of labor productivity to land productivity. In simplified notation, the central condition is that
\[
  \mathbb{E}_0\sum_t\left(\frac{A_{Ht}}{A_{Xt}}\right)^{1/\sigma-1}
\]
converges. Economically, \(\sigma>1\) makes the relevant exponent negative. Faster growth in labor productivity than land productivity suppresses land-rent growth, while rising income keeps pushing savings into land. Land prices rise faster than land rents can justify because the asset remains a savings vehicle even as its dividend becomes relatively weaker \citep[pp.~4, 13]{hiranoToda2025UnbalancedGrowth}.

\paragraph{P22.} Land overvaluation is not simply irrational exuberance. It is a model result in which the price of land separates from the present value of land rents under structural growth conditions. It also differs from a generic housing-price boom. Housing combines land and structures, so an observed housing rent-yield can remain stable even if the pure land component contains a bubble. Toda's measurement warning belongs to the same architecture: the analyst must separate the rent on structures from the pure land dividend before drawing conclusions about bubbles. Without that separation, evidence against a housing bubble can be misread as evidence against land overvaluation \citep[pp.~1--3]{toda2025LandBubbles}.

\paragraph{P23.} The leverage and asset-bubble branch extends overvaluation through a different mechanism. Hirano, Jinnai, and Toda model a financial accelerator in which land prices, capital investment, land rents, wealth, and asset demand interact. Rising land prices can raise aggregate wealth and collateral values, feeding investment and asset demand in a self-reinforcing process. This is not the crowding-out mechanism in the Hirano--Stiglitz credit papers. The crucial object is a leverage threshold. Below it, the economy can remain close to a fundamental path in which land prices reflect rent values. Above it, leverage strengthens the feedback from wealth to investment and from investment to land demand, producing a phase transition toward a nonstationary bubble path \citep[pp.~2--3, 9]{hiranoJinnaiToda2024Leverage}.

\paragraph{P24.} The synthesis is heterogeneous. One branch shows land speculation crowding out capital. Another shows leverage and land-price increases reinforcing capital investment up to a phase transition. That difference must remain visible. The effect of land-price inflation depends on the model architecture: whether land absorbs savings away from capital, whether collateral values relax investment constraints, whether leverage magnifies asset demand, and whether the equilibrium is fundamental or bubbly. A serious application to land-rent financialization cannot cite ``the literature'' as if it made one uniform claim.

\begin{center}
\small
\begin{longtable}{p{0.18\textwidth}p{0.22\textwidth}p{0.24\textwidth}p{0.18\textwidth}p{0.16\textwidth}}
\caption{Two mechanisms linking land valuation and capital formation}\label{tab:mechanisms-land-capital}\\
\toprule
Mechanism & Starting condition & Transmission channel & Endpoint & Caution\\
\midrule
\endfirsthead
\toprule
Mechanism & Starting condition & Transmission channel & Endpoint & Caution\\
\midrule
\endhead
Crowding-out & Finite savings or borrowing capacity; land competes with productive capital & Higher \(P_t\) absorbs savings, down payments, or collateral capacity & Lower \(k_{t+1}\), lower growth, possible instability & Do not read every land boom as crowding-out without the savings or credit-allocation condition.\\
Leverage accelerator & Productive agents can lever beyond a threshold & Land prices, wealth, collateral, investment, and asset demand reinforce one another & Transition from fundamental to bubbly growth path & Do not collapse this into the crowding-out branch; here leverage can initially raise investment.\\
\bottomrule
\end{longtable}
\end{center}

\section{Mechanism IV: \texorpdfstring{\(R\)}{R} versus \texorpdfstring{\(G\)}{G}, land, and infinite debt rollover}

\paragraph{P25.} The \(R\)-versus-\(G\) branch starts from a conventional result. In overlapping-generations economies with land, McCallum and Homburg are read as showing that a productive, non-reproducible, non-depreciating asset such as land rules out the landless Diamond-type oversaving problem. Along a balanced growth path, land rents and prices grow with output. Because land pays a positive rent, the return to holding land combines rent with price appreciation. If land prices grow with output, that rent component makes the total land return exceed the growth rate. By no arbitrage, the safe interest factor is pushed above the growth factor, \(R>G\), and infinite debt rollover is ruled out \citep[pp.~1--2, 16--17]{hiranoToda2025LandInfiniteDebtRollover}.

\paragraph{P26.} Hirano and Toda narrow that conventional land-efficiency result. They argue that it depends on knife-edge balanced-growth restrictions, including equal productivity growth or Cobb--Douglas-like substitution conditions. Once the model allows unbalanced growth, with labor productivity growing faster than land productivity and \(\sigma>1\), land rents need not grow at the same rate as output. Land can become asymptotically negligible relative to aggregate endowments even though it remains an asset held for savings. Under those broader conditions, a Pareto-inefficient equilibrium with \(R<G\) can arise even in an economy with land, reopening the possibility of infinite debt rollover \citep[pp.~1--2, 16--17]{hiranoToda2025LandInfiniteDebtRollover}.

\paragraph{P27.} The 2026 CIGS version cautions against using \(R<G\) as a policy shortcut. The paper criticizes the heuristic that market frictions generating \(R<G\) automatically settle the debt-rollover question. If \(G\), the relevant asset return, and the market frictions are not specified, the inequality loses theoretical coherence \citep[p.~4]{hiranoToda2026LandGversusR}. The inequality is not self-interpreting. It depends on which return is being compared with which growth rate, whether land is priced by the safe rate or by a constrained asset return, and whether the equilibrium is balanced, unbalanced, fundamental, or bubbly. This matters for the Fondecyt project because Latin American debt and land markets are open, institutionally heterogeneous, and exposed to currency, sovereign-risk, and external-finance conditions outside the closed OLG benchmark. \projecttranslation{}

\paragraph{P28.} The debt-rollover result should not carry more weight than the model gives it. The formal result concerns infinite debt rollover in OLG economies with land and unbalanced growth. It does not by itself theorize private developer balance sheets, mortgage refinancing, or corporate debt rollover in Latin American urban land markets. The useful conceptual move is narrower: land can absorb savings, carry bubble value, and alter the relation between asset returns and growth. Translating that into owner, developer, or household refinancing requires a different institutional account of liabilities, maturities, liquidity pressure, and credit-market access. The executive memo's translation, from OLG sellers to financialized owners facing rollover needs, is therefore a project-level extension rather than a result already proven in the SHT corpus. \projecttranslation{}

\section{Wobbly dynamics and endogenous instability}

\paragraph{P29.} The wobbly-dynamics papers explain why land-price dynamics need not converge smoothly. In the global-dynamics branch, rational-expectations economies can fluctuate indefinitely without settling into a regular steady state. The underlying object is a multiplicity of momentary equilibria: for a given current state, different expectation-consistent paths can clear the market. Land widens the parameter region in which such dynamics occur because it offers an alternative store of value that competes with capital and changes the relation among savings, capital accumulation, interest rates, and asset prices \citep[pp.~1--2, 10]{hiranoStiglitz2022WobblyEconomy}; \citep[pp.~7, 24, 46--47]{hiranoStiglitz2022LandSpeculation}.

\paragraph{P30.} The branch contributes a theory of endogenous instability. Instability is not imposed from outside by irrational shocks. It follows from rational expectations, finite life, asset-market interactions, and phase transitions. A land-price boom can remain consistent with expectations for a time. As land absorbs savings, capital falls, wages and output fall, interest rates rise, and the required land-price appreciation becomes harder to sustain. The threshold is endogenous: it is generated by the same arbitrage and savings relations that initially sustain the boom. At that point, the economy shifts to a crash path \citep[pp.~7, 24, 46--47]{hiranoStiglitz2022LandSpeculation}.

\paragraph{P31.} Wobbly dynamics remains supplementary to the review's main argument. It gives the project a language for instability, bounded boom-bust paths, and endogenous crashes, but it does not replace the credit-allocation, overvaluation, or \(R\)-versus-\(G\) mechanisms. Its value is to prevent a static reading of land overvaluation. If land prices, credit, and capital formation interact dynamically, the question is not only whether land is overvalued at one date. The question is how valuation paths affect investment, collateral, refinancing, and crisis timing. The specific translation to Latin American boom-bust paths remains a \sourcegap{}.

\section{Policy implications inside the corpus}

\paragraph{P32.} The policy implication is not less credit in the abstract. Composition matters. Credit to productive capital can crowd in investment and growth; credit to real estate or land can crowd out capital and retard growth. A policy agenda drawn from the corpus would distinguish productive collateral rules from land-collateral rules, then ask whether financial liberalization raises capital formation or mainly increases leveraged land demand. In the notation used by the credit models, the contrast is between easing \(\theta\), which raises the pledgeability of productive capital, and easing \(\theta^x\), which raises the pledgeability of land. The policy target is therefore the channel that turns liquidity into land demand, not credit creation as such \citep[pp.~1, 15, 33, 38]{hiranoStiglitz2024CreditLandGrowth}; \citep[p.~1]{hiranoStiglitz2025LowInterestPolicy}.

\paragraph{P33.} Land taxation responds directly to the speculative allocation problem. In the Henry George paper, land taxation is not only a distributional device. It changes the return to land speculation and can redirect savings toward productive capital. The corpus also points to related instruments: land-rent taxation, capital-gains taxation, public investment financed through land rents, and collateral regulation that limits the use of land as a speculative borrowing base. That policy claim remains model-dependent: it assumes that the relevant distortion is land speculation and that the state can tax land or land rents without reproducing other distortions \citep[pp.~327--328, 335, 344]{hiranoStiglitz2025HenryGeorge}. In Latin American settings, cadastres, informality, municipal fiscal capacity, political capture, and legal property regimes would have to be incorporated before the policy can be operationalized. \sourcegap{}

\paragraph{P34.} Monetary policy has the same compositional limit. Lower interest rates can support productive investment if the financial system channels liquidity toward capital formation. They can also raise the leveraged return to land and real estate. The SHT contribution is to show that low rates are not automatically pro-growth when land is a major store of value and collateral object. The relevant question is whether lower rates reduce the cost of productive investment or instead lower the financing cost of acquiring a fixed asset whose rising price then crowds out capital. The policy implication is targeted regulation: restrict the channels that convert liquidity into land-price inflation while preserving or expanding credit for productive investment \citep[pp.~1--4]{hiranoStiglitz2025LowInterestPolicy}; \citep[pp.~1, 15, 33, 38]{hiranoStiglitz2024CreditLandGrowth}.

\paragraph{P35.} Measurement is a policy issue as well. If wealth statistics rise because land values or capitalized rents rise, policy makers may mistake balance-sheet expansion for improved sustainability. Stiglitz's wealth measurement work warns against that interpretation. Toda's housing-rent warning adds a second problem: the analyst must separate land from structures before interpreting rent yields or price-rent ratios. A land-rent financialization agenda should therefore track asset prices and wealth-to-income ratios alongside productive capital formation, sectoral credit allocation, land-rent capitalization, and the separation between structures and land \citep[pp.~1--2, 4, 11]{stiglitz2015MeasurementWealth}; \citep[pp.~1--3]{toda2025LandBubbles}.

\section{Critical appraisal and Fondecyt translation}

\paragraph{P36.} The corpus is strongest as conceptual and mechanistic theory. It lets the project argue that land-price inflation can be macroeconomically harmful without relying on a moral claim about speculation. Land is fixed and non-produced. Land prices can capitalize rents. Land collateral can expand credit. Credit can move toward or away from productive capital. Unbalanced growth can separate land prices from rents. Debt-rollover conditions depend on the asset structure of the economy. These mechanisms give the Fondecyt project a disciplined theoretical architecture.

\paragraph{P37.} Transportability is the weak point. The models are stylized: OLG structures, rational expectations, closed or simplified economies, representative sectors, and abstract land markets. Their limits fall into four groups. First, the model form abstracts from historically specific classes, firms, and state institutions. Second, land appears as a homogeneous asset rather than legally and spatially differentiated urban property. Third, the credit system is formalized through collateral parameters rather than segmented banking, informal credit, or state-directed finance. Fourth, urban spatial structure appears only indirectly, if at all. These models do not directly explain Latin American urban land tenure, informal settlement, developer finance, state planning, zoning, mortgage segmentation, dollarized property markets, extractive land rents, or peripheral financial dependence. The omissions do not invalidate the mechanisms. They define the work required before those mechanisms become claims about Latin America. \projecttranslation{} \sourcegap{}

\paragraph{P38.} The executive memo's first major translation concerns OLG. In the SHT branch, OLG solves a formal asset-circulation problem: young agents buy, old agents sell, and resale expectations can support land demand. The Fondecyt project need not adopt OLG as its ontology. It can translate the same circulation function into a political-economy account of owners, developers, households, or firms that sell land or real estate because they face liabilities, maturities, liquidity constraints, refinancing pressure, and balance-sheet management. The translated mechanism would no longer be ``young buy, old sell''; it would be ``asset holders refinance, roll over, or liquidate claims under financial pressure.'' The corpus has not established that translation. It should be developed as a project contribution and marked as such. \projecttranslation{} \sourcegap{}

\paragraph{P39.} A second translation concerns unbalanced growth. In Hirano and Toda, unbalanced growth is a formal condition about differential productivity growth and substitution between land and non-land factors. In the project's developing framework, unbalanced growth may refer to a non-proportional relation among demand, capital accumulation, credit, land valuation, and productive capacity. These are not the same concept.

\begin{center}
\small
\begin{longtable}{p{0.19\textwidth}p{0.28\textwidth}p{0.28\textwidth}p{0.17\textwidth}}
\caption{Two uses of ``unbalanced growth''}\label{tab:unbalanced-growth-uses}\\
\toprule
Concept & SHT meaning & Project-facing meaning & Drafting rule\\
\midrule
\endfirsthead
\toprule
Concept & SHT meaning & Project-facing meaning & Drafting rule\\
\midrule
\endhead
Unbalanced growth in SHT & Differential productivity growth, usually labor or human-capital productivity growing faster than land productivity, with substitution conditions such as \(\sigma>1\). & Formal condition that helps produce land overvaluation or reopen debt-rollover results. & Use only when reconstructing the theorem or debt-rollover models.\\
Broader disproportional growth & Non-proportional movement among demand, capital accumulation, credit, land valuation, and productive capacity. & A Fondecyt concept still to be developed. & Mark as project translation and do not attribute it to SHT.\\
\bottomrule
\end{longtable}
\end{center}

The review should use the SHT concept when reconstructing the theorem and reserve the project's broader disproportional-growth concept for later theoretical development. Collapsing the two would create false continuity. \projecttranslation{}

\paragraph{P40.} The corpus leaves empirical work open. It motivates questions about sectoral credit, real-estate collateral, land taxation, and land-price/rent separation, but it does not answer them for Chile or Latin America. Some data needs follow directly from SHT variables: sectoral credit flows, collateral rules, land prices, land rents, produced capital, interest rates, and growth rates. Other data needs belong to the translation layer: land and structure decomposition, mortgage and developer finance, property taxation, informality, zoning, tenure, metropolitan price gradients, and municipal fiscal capacity. Until that work is done, the review can identify mechanisms and research hypotheses, not validated regional facts. \projecttranslation{} \sourcegap{}

\section{Conclusion}

\paragraph{P41.} The \SHT{} branch is best read as a theory of how a non-produced asset can reorganize macro-financial relations. Land is not just another input. It is a store of value and collateral object whose price can rise without expanding productive capacity. That price increase can absorb savings, redirect credit, change collateral values, alter growth paths, and affect the relation between returns and growth. The branch gives the Fondecyt project a rigorous way to connect land rent, credit, capital formation, overvaluation, and debt rollover.

\paragraph{P42.} The most defensible mechanisms are the wealth-capital distinction, the sectoral-credit distinction, the savings/capital crowding-out channel, the unbalanced-growth route to land overvaluation, the measurement distinction between housing rent and pure land rent, and the \(R\)-versus-\(G\) critique of conventional land-efficiency results. They are defensible because the corpus supplies formal mechanisms and source-routed page anchors. The fragile elements are the direct application of OLG to Latin American owners, the leap from public infinite debt rollover to private balance-sheet refinancing, and any empirical claim about regional land-rent financialization not yet supported by regional data. Those elements should enter the next stage as research questions, not as established findings. \projecttranslation{} \sourcegap{}

\paragraph{P43.} The next phase should keep the order disciplined: define the objects, reconstruct the mechanisms, state the causal chains, identify variables and parameters, distinguish theory from evidence, and only then translate. The review should not claim that SHT already explains Latin American land-rent financialization. It should claim that SHT provides a formal mechanism set that can be translated, tested, and modified for that project.

\section*{Citation and source-status notes}

\begin{itemize}
  \item \sourcegap{} and \projecttranslation{} mark claims that belong to the Fondecyt translation layer rather than the bounded SHT corpus.
  \item Internal note layers were not cited as sources. They were used only to locate claims and page anchors in the original papers.
  \item Gasic (2026) was used only as a project-facing guardrail for corpus separation. It is not cited as support for SHT mechanisms, theorems, empirical motivations, or policy results.
\end{itemize}

% ---------- Bibliography ----------%
\pagebreak
\bibliographystyle{apalike}
\bibliography{references}

\end{document}
